% =====================================================================
% บทที่ 3 การเริ่มต้นใช้งาน (ผู้อ่าน: ผู้ใช้ทุกบทบาท)
%
% วิธีวางหมายเลขบนภาพ:
%   1) ใส่ตัวเลือก grid ชั่วคราว  \begin{screenshot}[grid]{images/login.png}{...}{...}
%   2) อ่านพิกัดจากเส้นตาราง (0 = ซ้าย/บน, 1 = ขวา/ล่าง) แล้วใส่ \callout{x}{y}{เลข}
%   3) ลบ grid ออกก่อนส่งเล่ม
% =====================================================================
\label{ch:start}

\begin{guide}
คู่มือผู้ใช้งานทั่วไปเริ่มที่บทนี้ อธิบายงานแรกที่ผู้ใช้ทุกบทบาทต้องทำ ได้แก่ สมัครสมาชิก เข้าสู่ระบบ
รู้จักหน้าจอหลัก และออกจากระบบ
\begin{itemize}
  \item ภาพหน้าจอหลักพร้อมหมายเลขกำกับ และคำอธิบายทุกหมายเลขใต้ภาพ
  \item เงื่อนไขที่ผู้ใช้มักพลาด เช่น รูปแบบอีเมล ความยาวรหัสผ่าน ใส่ในกล่องสำคัญ (important)
  \item ภาพที่กว้างแต่เตี้ย เช่น แถบเมนู ให้ครอปเฉพาะส่วนแล้ววางหมายเลขนอกภาพ
\end{itemize}
แอปบนมือถือ: เพิ่มการเปิดแอปครั้งแรก การอนุญาตสิทธิ์ (กล้อง ตำแหน่ง การแจ้งเตือน) และท่าทางสัมผัสที่ใช้
\end{guide}

\section{การสมัครสมาชิก}
\label{sec:register}

ผู้ใช้ใหม่ทุกบทบาทต้องสมัครสมาชิกก่อนเข้าใช้งาน โดยใช้อีเมลของมหาวิทยาลัยเท่านั้น

\begin{steps}
  \item เปิดหน้าเข้าสู่ระบบ แล้วคลิกลิงก์ \btn{สมัครสมาชิก} ใต้ปุ่มเข้าสู่ระบบ
  \item กรอกข้อมูลในแบบฟอร์มให้ครบทุกช่อง
        \begin{enumerate}
          \item \field{อีเมล} และ \field{รหัสผ่าน}
          \item \field{ชื่อ-นามสกุล (ภาษาไทย)} และ \field{ชื่อ-นามสกุล (ภาษาอังกฤษ)}
          \item \field{บทบาท} ของผู้ใช้
        \end{enumerate}
  \item ถ้าเลือกบทบาทนักศึกษา ให้กรอก \field{รหัสนักศึกษา} เช่น \cmd{663440001}
  \item คลิก \btn{สมัครสมาชิก}
        \result{ระบบสร้างบัญชีและพากลับไปยังหน้าเข้าสู่ระบบ}
\end{steps}

\begin{important}
อีเมลต้องลงท้ายด้วย \cmd{@kku.ac.th} เท่านั้น และรหัสผ่านต้องมีอย่างน้อย 6 ตัวอักษร
ระบบไม่รับอีเมลจากผู้ให้บริการอื่น เช่น Gmail
\end{important}

\section{การเข้าสู่ระบบ}
\label{sec:login}

เปิดเว็บเบราว์เซอร์ไปที่ \systemurl{} จะพบหน้าเข้าสู่ระบบดังภาพที่~\ref{fig:login}
ซึ่งประกอบด้วยส่วนต่าง ๆ ตามหมายเลขกำกับ

\begin{screenshot}[width=0.52\linewidth]{images/login.png}{หน้าเข้าสู่ระบบ}{fig:login}
  \callout{0.94}{0.49}{1}
  \callout{0.94}{0.65}{2}
  \calloutline{0.845}{0.555}{0.845}{0.625}{3}
  \highlight{0.10}{0.735}{0.80}{0.095}
  \callout{0.94}{0.78}{4}
  \callout{0.79}{0.873}{5}
\end{screenshot}

\begin{calloutlist}
  \item ช่องกรอกอีเมลของมหาวิทยาลัย
  \item ช่องกรอกรหัสผ่าน
  \item ปุ่มแสดงหรือซ่อนรหัสผ่านที่พิมพ์
  \item ปุ่มเข้าสู่ระบบ
  \item ลิงก์ไปยังหน้าสมัครสมาชิก (หัวข้อ~\ref{sec:register})
\end{calloutlist}

\begin{steps}
  \item กรอกอีเมลในช่องหมายเลข~\calloutmark{1} และรหัสผ่านในช่องหมายเลข~\calloutmark{2}
  \item คลิก \btn{เข้าสู่ระบบ} หรือกด \keys{Enter}
        \result{ระบบแสดงหน้ารายการโครงงานตามบทบาทของผู้ใช้ (หัวข้อ~\ref{sec:main-screen})}
\end{steps}

\begin{note}
หากลืมรหัสผ่าน ให้ติดต่อผู้ดูแลระบบเพื่อรีเซ็ตรหัสผ่าน (หัวข้อ~\ref{sec:contact})
\end{note}

\section{ส่วนประกอบของหน้าจอหลัก}
\label{sec:main-screen}

หลังเข้าสู่ระบบ ผู้ใช้จะเห็นหน้ารายการโครงงาน ซึ่งแบ่งเป็น 2 ส่วน ได้แก่
แถบเมนูด้านบน (ภาพที่~\ref{fig:topbar}) ที่ใช้ได้จากทุกหน้าของระบบ
และพื้นที่รายการโครงงาน (ภาพที่~\ref{fig:dashboard})

% ภาพที่กว้างแต่เตี้ย เช่น แถบเมนู: ครอปเฉพาะส่วนนั้นให้ตัวอักษรอ่านออก
% แล้ววางหมายเลขไว้ "นอกภาพ" ด้วยค่า y มากกว่า 1 พร้อมเส้นชี้ขึ้นไปยังปุ่ม
\begin{screenshot}[width=0.65\linewidth]{images/dashboard-topbar.png}{แถบเมนูด้านบน}{fig:topbar}
  \calloutline{0.088}{1.7}{0.088}{0.82}{1}
  \calloutline{0.287}{1.7}{0.287}{0.82}{2}
  \calloutline{0.527}{1.7}{0.527}{0.82}{3}
  \calloutline{0.678}{1.7}{0.678}{0.82}{4}
  \calloutline{0.843}{1.7}{0.843}{0.82}{5}
\end{screenshot}

\begin{calloutlist}
  \item การแจ้งเตือน เช่น ข้อเสนอได้รับการอนุมัติหรือมีความเห็นใหม่
  \item ชื่อผู้ใช้และบทบาทที่เข้าสู่ระบบอยู่
  \item เมนู \btn{ส่งงาน} สำหรับส่งเอกสารตามกำหนด (หัวข้อ~\ref{sec:submit})
  \item ปุ่มตั้งค่าบัญชีผู้ใช้
  \item ปุ่ม \btn{ออกจากระบบ}
\end{calloutlist}

\begin{screenshot}{images/dashboard.png}{พื้นที่รายการโครงงาน}{fig:dashboard}
  \callout{0.625}{0.155}{1}
  \callout{0.968}{0.155}{2}
  \callout{0.38}{0.314}{3}
  \calloutline{0.40}{0.51}{0.313}{0.51}{4}
  \calloutline{0.40}{0.67}{0.331}{0.67}{5}
  \calloutline{0.40}{0.83}{0.144}{0.83}{6}
\end{screenshot}

\begin{calloutlist}
  \item ช่องค้นหาโครงงานจากชื่อ
  \item ปุ่ม \btn{+ สร้างข้อเสนอ} (หัวข้อ~\ref{sec:proposal})
  \item แท็บกรองโครงงานตามสถานะ ได้แก่ ทั้งหมด ฉบับร่าง รอพิจารณา อนุมัติ และแก้ไข
  \item สถานะล่าสุดของโครงงาน
  \item การ์ดโครงงาน แสดงชื่อโครงงาน นักศึกษา และอาจารย์ที่ปรึกษา
  \item ปุ่ม \btn{ดูรายละเอียด} เปิดหน้ารายละเอียดของโครงงาน
\end{calloutlist}

\section{การออกจากระบบ}

คลิก \btn{ออกจากระบบ} ที่มุมขวาบน (หมายเลข~\calloutmark{5} ในภาพที่~\ref{fig:topbar})
ระบบบันทึกเอกสารอัตโนมัติทุก 2 วินาทีที่มีการแก้ไข แต่ควรกด \keys{Ctrl+S} ก่อนออกจากหน้าแก้ไขเอกสารทุกครั้ง
